\documentclass[a4paper, 10pt]{article}
% CUERPO
\usepackage[margin=2.5cm]{geometry}
\usepackage{graphicx}
\graphicspath{{images/}}
\begin{document}
\section{Analisis Cualitativo}
\subsection{Clasificar los 16 estilos románticos}
\ El análisis del comportamiento por signo se basa en la fórmula hallada.
 Usando como base la fórmula: 
$$ k1,k2 = \frac{(a+d) \pm \sqrt{(a+d)^2 -4(ad - bc)} }{2}$$
se desprende el criterio traza-determinante:
\begin{enumerate}
\item Discriminante menor que 0 =$>$ Soluciones complejas (Foco)
\begin{enumerate}
\item (a+d) $>$ 0 =$>$ Inestable
\item (a+d) $<$ 0 =$>$ Estable
\end{enumerate}
\item Discriminante mayor que 0 =$>$ Soluciones reales
\begin{enumerate}
 \item $|TrA|>|\sqrt{Disc}|$ =$>$ Nodo (Mismo signo)
  \begin{enumerate}
    \item $k1<0$ =$>$ Estable
    \item $k1>0$ =$>$ Inestable
  \end{enumerate}
   \item $|TrA|<|\sqrt{Disc}|$ =$>$ Silla(Inestable)(Diferente signo)
  \end{enumerate}
\end{enumerate}
Traza-Determinante:
\begin{enumerate}
    \item $det(A)<0$ =$>$ Silla
    \item $det(A)>0$ y $\Delta<0$  =$>$ Foco
    \item $det(A)>0$ y $\Delta>0$  =$>$ Nodo
    \item $det(A)>0 $ y $\Delta=0$ =$>$ Nodo Degenerado
    \item $det(A)>0 $, $TrA=0$ y $\Delta=0$ =$>$ Centro
\end{enumerate}

\subsubsection{$a>0,\ b>0,\ c>0,\ d>0$} 
\ Si se cumple esta condicion, entonces la ecuacion resulta coincidente con el criterio; luego el análisis es análogo.
\subsubsection{$a>0,\ b>0,\ c>0,\ d<0$} 
\ Si el único negativo es d, luego la fórmula resulta en:
$$ k1,k2 = \frac{(a-d) \pm \sqrt{(a-d)^2 -4(-ad - bc)} }{2}$$
Luego el determinante es negativo, es una silla
\subsubsection{$a>0,\ b>0,\ c<0,\ d>0$}
Si c es menor que 0, entonces queda:
$$ k1,k2 = \frac{(a+d) \pm \sqrt{(a+d)^2 -4(ad + bc)} }{2}$$
Como el determinante es mayor que 0, nunca es una silla, por lo tanto depende del parametro, pudiendo ser nodo o foco
\subsubsection{$a>0,\ b>0,\ c<0,\ d<0$} 
Si c y d son menores que 0, entonces:
$$ k1,k2 = \frac{(a-d) \pm \sqrt{(a-d)^2 -4(-ad +bc)} }{2}$$
Este caso depende del valor de los parametros, luego se ajusta al criterio inicial
\subsubsection{$a>0,\ b<0,\ c>0,\ d>0$}
Si b es negativo tenemos:
$$ k1,k2 = \frac{(a+d) \pm \sqrt{(a+d)^2 -4(ad + bc)} }{2}$$
Luego como el determinante es positivo, nunca es una silla, el analisis depende del parametro, pudiendo ser nodo o foco
\subsubsection{$a>0,\ b<0,\ c>0,\ d<0$}
Si b y d son menores que 0, entonces la formula es:
$$ k1,k2 = \frac{(a-d) \pm \sqrt{(a-d)^2 -4(-ad + bc)} }{2}$$
Esta formula es analoga a c y d menores que 0
\subsubsection{$a>0,\ b<0,\ c<0,\ d>0$} 
Si b y c son negativos, entonces:
 $$ k1,k2 = \frac{(a+d) \pm \sqrt{(a+d)^2 -4(ad - bc)}}{2}$$
Donde  $-(-bc)=bc$, por lo tanto es la formula original, por  lo tanto me ajusto al criterio
\subsubsection{$a>0,\ b<0,\ c<0,\ d<0$} 
Si b,c y d son negativos, entonces:
$$ k1,k2 = \frac{(a-d) \pm \sqrt{(a-d)^2 -4(-ad - bc)}}{2}$$
El determinante siempre es negativo, por lo tanto es silla
\subsubsection{$a<0,\ b>0,\ c>0,\ d>0$} 
Si a es negativo, entonces:
$$ k1,k2 = \frac{(-a+d) \pm \sqrt{(-a+d)^2 -4(-ad - bc)} }{2}$$
El proceso es analogo al anterior, luego es una silla, el determinante es negativo
\subsubsection{$a<0,\ b>0,\ c>0,\ d<0$} 
Si a y d son negativos, entonces:
$$ k1,k2 = \frac{(-a-d) \pm \sqrt{(-a-d)^2 -4(ad - bc)} }{2}$$
En este caso no podemos decir nada, puesto que depende de los parámetros
\subsubsection{$a<0,\ b>0,\ c<0,\ d>0$} 
Si a y c son negativos, entonces:
$$ k1,k2 = \frac{(-a+d) \pm \sqrt{(-a+d)^2 -4(-ad + bc)} }{2}$$
No tengo que decir en este caso, depende del valor entre los parametros, luego me ajusto al criterio
\subsubsection{$a<0,\ b>0,\ c<0,\ d<0$} 
Si a, c y d son negativos, luego:
$$ k1,k2 = \frac{(-a-d) \pm \sqrt{(-a-d)^2 -4(ad + bc)} }{2}$$
Este caso depende del valor de los parametros, como el determinante es positivo,
entonces puede ser un foco o un nodo
\subsubsection{$a<0,\ b<0,\ c>0,\ d>0$} 
Si a y b son negativos, entonces:
$$ k1,k2 = \frac{(-a+d) \pm \sqrt{(-a+d)^2 -4(-ad + bc)} }{2}$$
El problema depende de los parámetros, luego se ajusta al criterio inicial
\subsubsection{$a<0,\ b<0,\ c>0,\ d<0$} 
Si a,b y d son negativos, luego:
$$ k1,k2 = \frac{(-a-d) \pm \sqrt{(-a-d)^2 -4(ad + bc)} }{2}$$
El caso depende del valor de los parametros, como el determinante es positivo,
pueder ser una nodo o un foco.
\subsubsection{$a<0,\ b<0,\ c<0,\ d>0$} 
Si a,b,c son negativos, entonces:
$$ k1,k2 = \frac{(-a+d) \pm \sqrt{(-a+d)^2 -4(-ad - bc)} }{2}$$
Como el determinante siempre es negativo, luego es una silla, por el criterio traza-determinante.
\subsubsection{$a<0,\ b<0,\ c<0,\ d<0$}
Si todos son negativos, entonces:
$$ k1,k2 = \frac{(-a-d) \pm \sqrt{(-a-d)^2 -4(ad - bc)} }{2}$$
Este caso depende del valor de los parametros, luego se ajusta al criterio inicial


\subsection{Dibujar retratos de fase para cada caso }
\ El analisis de los retratos de fase para cada combinación de simbolos puede resumirse en el analisis
del comportamiento general de los tipos de solución.
\subsubsection{Silla}
\ Las sillas son un tipo de solución que siempre son inestables, sus valores
propios son reales, y como presentan diferentes signos entonces sus soluciones se
acercan al origen y luego se alejan del origen.
\begin{figure}[h]
    \centering
    \includegraphics[width=0.5\linewidth]{Silla.png}
    \caption{Sillas}
    \label{fig:ejemplo}
\end{figure}
\subsubsection{Nodo}
\ Los nodos son un tipo de solución que pueder ser inestables o estables, sus valores 
propios son reales, y son ambos presentan el mismo signo; de aquí se desprenden los casos:
\begin{enumerate}
  \item $\lambda 1 <0$: Todas las soluciones tienden al origen
  \item $\lambda 1 >0$: Todas las soluciones se alejan del origen
\end{enumerate}
\begin{figure}[h]
    \centering
  \begin{minipage}[b]{0.45\textwidth}
    \includegraphics[width=\textwidth]{NodoEstable.png}
    \caption{Nodo Estable}
  \end{minipage}
  \hfill
  \begin{minipage}[b]{0.45\textwidth}
    \includegraphics[width=\textwidth]{NodoInestable.png}
    \caption{Nodo Inestable}
  \end{minipage}
\end{figure}
\subsubsection{Foco}
\ Los focos son un tipo de solución que aparecen cuando las soluciones
del sistema se vuelven complejas, osea, sus eigenvalores son complejos. 
La solución general de estos sistemas es de la forma:
\begin{equation}
Y_2(t)=e^{\alpha t}
(B_1\sin(\beta t)+B_2\cos(\beta t)).
\end{equation}
En la ecuación:\\
\ $\alpha$:  determina si crece o decrece la amplitud \\
\ $\beta$ : determina la rapidez de giro, su periodo natural es
$\frac{2\pi}{\beta}$.\\
Estos poseen varios subcasos, si tomamos que: $\lambda = \alpha \pm \beta i$:
\begin{enumerate}
  \item Estables: Si $\alpha<0$ (Las trayectorias se acercan girando al origen)
  \item Inestables: Si $\alpha>0$ (Las trayectorias se alejan girando del origen) 
  \item Centro: Si $\alpha=0$ (Curvas cerradas periódicas)
\end{enumerate}
\begin{figure}[h]
      \centering
  \begin{minipage}[b]{0.45\textwidth}
    \includegraphics[width=\textwidth]{FocoEstable.png}
    \caption{Foco Estable}
  \end{minipage}
  \hfill
  \begin{minipage}[b]{0.45\textwidth}
    \includegraphics[width=\textwidth]{FocoInestable.png}
    \caption{Foco Inestable}
  \end{minipage}
   \hfill
  \begin{minipage}[b]{0.45\textwidth}
    \includegraphics[width=\textwidth]{Centro.png}
    \caption{Centro}
  \end{minipage}
\end{figure}

\subsection{Modelo no lineal}
\ El modelo no lineal, consiste en la hipótesis de: si el amor de Julieta es
abundante, Romeo responde de manera adversa, y si Julieta odia a Romeo, entonces
 Romeo decide tratarla mejor.
\ Este modelo puede expresarse como un modelo lógistico:
$$\frac{dR}{dt}=aR+bJ(1-|J|) $$
$$\frac{dJ}{dt}=cR(1-|R|)+dJ$$
\subsubsection{Puntos de Equilibrio}
Luego, los puntos de equilibrio son las soluciones del sistema igualado a cero, 
del cual:
$$0=aR+bJ(1-|J|) $$
$$0=cR(1-|R|)+dJ$$
Luego:
$$J=\frac{-cR(1-|R|)}{d}$$
Sustituyendo en $0=aR+bJ(1-|J|)$:
$$0=aR+b(\frac{-cR(1-|R|)}{d})(1-|\frac{-cR(1-|R|)}{d}|) $$
$$0=R(a+b(\frac{-c(1-|R|)}{d})(1-|\frac{-cR(1-|R|)}{d}|) ) $$
Fijando los valores a=b=-2 y c=d=1; obtenemos la ecuación:
$$0=R(-2 -2(-(1-|R|))(1-|-R(1-|R|)|) ) $$
El análisis de esta ecuación se debe realizar por los valores de R, debido al 
del módulo; se obtienen los equilibrios:
\begin{enumerate}
  \item (0,0)
  \item (2,2)
  \item (-2,-2)
\end{enumerate}
\subsubsection{Proceso de linealización}
 El proceso de linealización se basa en la idea de que, a partir de una vecindad del equilibrio,
  la función tiene un comportamiento similar a su derivada. Partiendo del (0,0):
 \\ \ Digamos: $R=0+u$ y $J=0+v$ :
 $$\frac{du}{dt} = -2u - 2v(1 - |v|)$$
 $$\frac{dv}{dt} = u(1-|u|) + v$$
 Luego como el módulo de u y v son cercanos a cero; podemos decir que
 $1-|u|$ y $1-|v|$ son cercanos a 1, luego:
 $$\frac{du}{dt} = -2u - 2v$$
 $$\frac{dv}{dt} = u + v$$
 El proceso es analogo para cada equilibrio, este se resume en:
 \begin{enumerate}
  \item Declarar $R = \epsilon x +u$ y  $J = \epsilon y +v$
  \item Sustituir en la ecuación los valores
  \item Analizar comportamiento
  \item Escribir el sistema matricial lineal
 \end{enumerate}
 \subsubsection{Ausencia de ciclos}
 La ausencia o presencia de ciclos en el sistema determinarse por el criterio de Bendixson:
 \\ Si se cumple que:
 $$\frac{dx}{dt} =f(x,y) $$
 $$\frac{dy}{dt} =g(x,y) $$
 $$ $$
 Y en una región simple conexa D:
 $$ \frac{\partial f}{\partial x} + \frac{\partial g}{ \partial y}$$
 
 Mantiene el signo constante y no es identicamente a cero, entonces no hay órbitas periódicas en D.
 \\ Como el sistema es de la forma:
$$\frac{dR}{dt}=aR+bJ(1-|J|) $$
$$\frac{dJ}{dt}=cR(1-|R|)+dJ$$
 Entonces:
  $$ \frac{\partial f}{\partial R} + \frac{\partial g}{ \partial J} =  a + d$$
  la existencia de ciclos depende de los valores de los parametros, como en el caso a=-2 y d=1, entonces es
  constante de valor igual a -1, por lo tanto no presenta ciclos límites. Hay que destacar que depende
   de la región del plano, puesto que esto afecta si R o J son positivos.
\end{document}